\chapter{Limitations and Threats to Validity}\label{ch:limitations}

This chapter is written to be read before the conclusions rather than after them. Each limitation
states what was not established, why, and what would be required to establish it. Nothing here is
hedging language around a result claimed elsewhere; where a limitation contradicts a stronger
reading of a result, the limitation is correct.

\section{The central limitation: no hardware measurement}

\textbf{No physical ESP32 was available in the environment that produced these results.} The
probe that established this is recorded in Appendix~\ref{app:hardware}: there were no serial
devices, no USB subsystem, no \code{arduino-cli}, no Python \code{serial} module, and the only
network interfaces present were \code{eth0}, \code{ifb0} and \code{lo}.

Four consequences follow, and they are the four most significant limitations of this work.

\textbf{No timing claim of any kind is supported.} The host harness advances a logical clock. The
safety latency is the same constant --- 151\,ms --- in every trial, by construction. Latency
variance in this environment is identically zero, which is not a finding about the firmware but a
property of the simulator.

\textbf{Logical isolation does not imply physical timing invariance.} This distinction is
maintained throughout the report and is restated here as a limitation rather than a nuance. The
model establishes that the safety task's enabling condition is independent of every network-side
variable, so no network behaviour can \emph{block} it. The firmware's asymmetric mutex timeout
makes the blocking component of the loop period zero by construction. Whether the \emph{total}
loop period is invariant to network load on real silicon is an empirical question about cache
contention, bus arbitration, interrupt latency, radio duty cycle and thermal behaviour, and it has
not been answered. Panel B of Figure~\ref{fig:setup} is labelled as not yet measured for this
reason and that label must not be removed until data exists.

\textbf{On-device memory safety is not established.} The three million fuzzing iterations ran
against a host build under AddressSanitizer and UndefinedBehaviorSanitizer. Neither sanitizer
exists for the Xtensa target. Compiler and architecture differences --- alignment requirements,
integer width, stack layout --- mean a host-clean parser is evidence for, but not proof of, an
Xtensa-clean parser.

\textbf{No radio-layer experiment was performed.} The deauthentication detector's parameters were
swept in simulation only. An RF-shielded enclosure is a hard prerequisite for a real
deauthentication experiment, both for validity and for legality, and none was available.

The pre-registered analysis plan for conditions A--D --- the Mann--Whitney U test, the sample
sizes, the acceptance criteria --- was committed \emph{before any data existed} and is reproduced
in Appendix~\ref{app:hardware}. Its purpose is that the test cannot be chosen after seeing the
measurements. The runbook in \code{docs/HARDWARE\_A\_TO\_C\_RUNBOOK.md} specifies the procedure in
enough detail to be executed by someone else.

\section{Threats to internal validity}

\subsection{The harness is not the firmware}

The firmware logic is compiled against stub headers in \code{tests/host/stubs/}. The logic itself
is unmodified --- the same \code{smart\_home.ino} is compiled, and only the headers it includes
differ --- but the environment around it is a model. There is no radio, no true parallel
execution, no contention for shared silicon and no real clock.

\subsection{Four defects were found in the measurement infrastructure itself}

This is reported as a threat to validity because two of the four produced results that flattered
the intended conclusion, and because it bears on how much confidence the remaining results
deserve.

\begin{enumerate}[leftmargin=1.8em,itemsep=4pt]
  \item \textbf{Unattributable availability.} The harness recorded that a command failed without
        recording why, so a refusal by policy and a refusal by rate limit were indistinguishable.
  \item \textbf{A token-bucket refill introduced by the first repair.} The repair to the above
        changed the arrival pattern, altering the quantity being measured.
  \item \textbf{A single-attempt measurement that measured bucket phase.} One attempt per trial
        measured where in the refill cycle the attempt landed rather than the admission rate.
        This is why the availability denominator is 1,250 rather than 300: multiple attempts per
        trial are required for the measurement to mean what it claims.
  \item \textbf{A 40\,ms button hold against a 40\,ms debounce.} A boundary condition that made
        the legitimate local path unreliable in the harness and had nothing to do with the
        firmware.
\end{enumerate}

A fifth defect, in the corpus rule, is described in Section~\ref{sec:res-corpus}: the rule
returned zero against this project's \emph{own} pre-repair guard, because it treated any mention
of the authentication identifier as evidence that authentication was required, and in the
defective guard that identifier appears negated.

Each repair is accompanied by a control that would fail if the defect recurred. That is the only
reason the remaining results are reported with confidence: not because the harness is assumed
correct, but because specific, named ways in which it was wrong are now tested against.

\subsection{The vacuous invariant}

For part of this project the formal specification passed every check while being structurally
incapable of failing (Section~\ref{sec:formal-vacuity}). A passing formal check was taken as
evidence for a property the check could not test. The repair, and the inverted configuration that
must report a violation, are described in Chapter~\ref{ch:formal}. The episode is a limitation of
the original method, not of the tool.

\subsection{No multiplicity correction}

Twelve ablation configurations are compared and two Fisher tests are reported. No correction for
multiple comparisons is applied, on the grounds that every arm is a fixed, pre-declared
configuration rather than a selected subset. This is stated so that a reader can judge the
decision rather than infer it. Both reported $p$-values are smaller than
$10^{-100}$ and would survive any correction, but that is a property of these particular results
and not a general defence of the choice.

\section{Threats to external validity}

\subsection{A single system}

This is a case study of one firmware artefact on one microcontroller with one broker. The defects
found are defects in this system. The \emph{categories} of defect --- an unverified configuration
assumption, a vacuous specification, an unbounded drain --- are argued to generalise; the specific
instances are not claimed to.

\subsection{The corpus is biased by construction}

The 135-project corpus was assembled from 22 fixed GitHub queries with a minimum of three stars
and an 80\,MB size cap, excluding archived repositories and forks. The criteria and rejection
counts are recorded in \code{corpus.json}.

GitHub projects with stars are not deployed products. The sample over-represents educational and
hobby code and under-represents commercial firmware, which is largely not public. \textbf{The
18-of-135 figure is not a prevalence estimate and must not be read as one}: rule 2
over-counts monitoring projects by roughly half, as the manual reading of all 18 showed, and
under-counts non-GPIO actuation by an unknown amount --- three of the nine confirmed cases actuate
through Modbus, RF or serial frames and were caught only by an incidental GPIO write.

Rule 1 returned \textbf{zero} hits. That is reported as a zero, and \textbf{no claim is made that
the guard pattern is widespread}.

\subsection{One broker, one version, one configuration}

Experiment E6 ran against Mosquitto 2.0.18 over TLS 1.2 with a private certificate authority.
Scenario E's containment depends on the access control list being deny-by-default and on the
device identity lacking write access on the command topic. A different broker, a different
version, or a permissive list would change the result. Other broker implementations ---
EMQX, HiveMQ, VerneMQ, AWS IoT Core --- were not tested, and their authorisation models differ.

\subsection{One framework, two versions}

Resource cost was measured under arduino-esp32 2.0.17 and 3.3.12, which differ by roughly 175\,kB
of flash. Behaviour under the ESP-IDF directly, or under other framework versions, was not
evaluated.

\section{Threats to construct validity}

\textbf{Alarm activation is not security.} In the \code{none}, \code{ratelimit} and \code{token}
arms the alarm fires in all 300 trials, and all three are completely defeated by
capture-and-replay. A single-metric evaluation would have concluded that no command defence
performs as well as a message authentication code. Three metrics are reported for this reason.

\textbf{Zero forgeries is not evidence of soundness.} Three million fuzz iterations producing no
accepted input is the expected outcome against a 128-bit tag whether the implementation is sound
or not, since a chance forgery has probability $2^{-128}$. The campaign is evidence of memory
safety and parser robustness, and is not independent evidence of authentication soundness.

\textbf{Simulated latency measures the simulator.} Stated above and repeated here because it is
the most likely misreading of Chapter~\ref{ch:results}.

\section{Scope exclusions}

\begin{itemize}[leftmargin=1.6em,itemsep=4pt]
  \item \textbf{Side-channel and fault-injection attack.} Out of scope; no equipment. Power
        analysis against the HMAC computation is a known attack class against embedded
        cryptography and was not evaluated.
  \item \textbf{Flash encryption and secure boot.} Not enabled. Section~\ref{sec:sec-flashdump}
        contains the flash-dump threat through authorisation rather than preventing extraction.
  \item \textbf{Over-the-air update security.} No update mechanism is implemented, so its
        security is neither claimed nor evaluated.
  \item \textbf{The controller host.} Treated as trusted infrastructure throughout. Scenario
        E$'$ is what happens when that assumption fails, and the layered build does not survive
        it.
  \item \textbf{Long-term operation.} No soak test. Flash wear from counter commits, clock drift,
        and memory fragmentation over months of operation were not evaluated.
  \item \textbf{Usability.} No human-factors study. The 41\% command-refusal rate under the token
        bucket is reported as a measured proportion; no claim is made about how users would
        respond to it.
\end{itemize}

\section{Known unresolved design problems}

Two problems were identified by this work and are \emph{not} fixed in the artefact. They are
recorded here rather than in future work, because a reader of the code should know they are open.

\textbf{Owner lockout under key compromise.} Two configurations admit 0 of 1,250 operator
commands, because the adversary advances the persisted replay counter beyond the operator's reach.
A candidate repair is a counter-resynchronisation path authenticated by the local button, which is
physically present and outside a remote adversary's reach. It has not been implemented or
evaluated, and the obvious alternatives --- a resynchronisation command over MQTT, or a counter
window --- weaken the replay resistance that the counter exists to provide.

\textbf{Scenario E$'$.} The layered configuration falls to a compromised controller host and
survives a device flash dump, which is the opposite of the cost ordering a defender would choose.
Mitigation requires an authority that does not reside on the controller --- a hardware security
module, a second factor for disarm, or a physically present confirmation --- and none is
implemented.

\section{Items requiring confirmation by the author or supervisor}

Recorded here so that no reader mistakes an unverified item for a verified one.

\begin{itemize}[leftmargin=1.6em,itemsep=4pt]
  \item The declaration, certificate and acknowledgement pages are unsigned drafts. The
        certificate page is explicitly marked as a template and carries no approval.
  \item \textbf{No supervisor approval has been obtained}, and no dissertation-overlap ruling has
        been made.
  \item \textbf{No institutional similarity screening report exists.} No similarity percentage is
        claimed anywhere in this report, official or otherwise.
  \item The competing-interest, funding and authorship-contribution declarations require explicit
        author confirmation before submission.
  \item No journal quartile is asserted anywhere in this work, because the ranking source could
        not be reached without circumventing bot protection, which was not attempted.
  \item Five references lack DOIs and require confirmation.
\end{itemize}
